\chapter*{Introduction Générale}
\mtcaddchapter

Aujourd’hui, les entreprises évoluent dans un environnement marqué par une transformation numérique rapide, les poussant à moderniser leurs méthodes de gestion afin d’améliorer leur efficacité et la qualité de leurs services. Dans ce contexte, la centralisation des données et l’automatisation des processus deviennent des éléments essentiels pour assurer un meilleur suivi des activités et une prise de décision plus efficace.

Le service après-vente (SAV), notamment dans le domaine de la maintenance technique, repose sur la gestion des clients, des équipements, des contrats et des interventions. Cependant, ces opérations sont souvent gérées de manière traditionnelle, ce qui entraîne des difficultés de suivi, un manque de coordination entre les acteurs et un risque d’erreurs ou de perte d’information.

Ce projet de fin d’études vise à concevoir et développer une application web de gestion du service après-vente dédiée à la maintenance des climatiseurs et des systèmes de surpression. L’objectif principal est de centraliser les informations, d’optimiser la gestion des interventions et d’améliorer la traçabilité des opérations ainsi que la communication entre les différents utilisateurs du système.

Afin de mener à bien ce projet, la méthodologie Scrum a été adoptée. Elle permet un développement itératif et progressif du système à travers plusieurs sprints, tout en garantissant une meilleure adaptation aux besoins et une amélioration continue du produit.

\begin{enumerate}
    \item Le \textbf{premier chapitre} présente l'étude préalable, le contexte du projet, ainsi que
          la problématique et la solution proposée.
    \item Le \textbf{deuxième chapitre} est consacré à la capture et la spécification des besoins.
    \item Le \textbf{troisième chapitre} traite le Sprint 1 : Authentification et référentiel métier.
    \item Le \textbf{quatrième chapitre} traite le Sprint 2 : Gestion des interventions.
    \item Le \textbf{cinquième chapitre} traite le Sprint 3 : Facturation et tableau de bord.
\end{enumerate}

\newpage
